% ============================================================
% Section 100: 特鲁德-洛伦兹介电函数与铜的光学性质
% ============================================================
\section{固体理论：特鲁德-洛伦兹介电函数与铜的光学性质}
\label{sec:drude-lorentz-dielectric}

\begin{modelbox}[题目：特鲁德-洛伦兹模型与复电导率]
固体介电常数的特鲁德-洛伦兹公式为
\begin{equation}
\varepsilon(\omega)=1+\frac{\omega_p^2}{(\omega_0^2-\omega^2)-i\omega\tau^{-1}},
\end{equation}
其中 $\omega_p$ 为等离子体频率，$\omega_0$ 为带间跃迁频率，$\tau$ 为电子散射时间。

(1) 室温下铜中 $\tau=10^{-14}\,$s，估计 $\omega_p$ 和 $\omega_0$ 的数量级。利用铜特有的``颜色''决定 $\omega_0$，定性画出 $\varepsilon(\omega)$ 的实部和虚部；
(2) 在室温下，计算铜的复电导率 $\sigma(\omega)$；
(3) 在绝对零度下，完全纯净无缺陷的铜的复电导率 $\sigma(\omega)$ 是什么？
\end{modelbox}

\begin{tipbox}[考点快速分析]
\begin{itemize}
    \item \textbf{等离子体频率}：$\omega_p=\sqrt{ne^2/(\varepsilon_0 m)}$
    \item \textbf{带间跃迁}：铜的红色对应蓝绿光被吸收，阈值约 $2\,$eV
    \item \textbf{介电函数与电导率关系}：$\varepsilon(\omega)=1+i\sigma(\omega)/(\varepsilon_0\omega)$
    \item \textbf{零温极限}：$\tau\to\infty$，耗散消失，电导率纯虚
\end{itemize}
\end{tipbox}

\begin{derivationbox}[标准解题过程]
\textbf{(1) $\omega_p$ 和 $\omega_0$ 的估计}

等离子体频率：铜的自由电子浓度 $n\approx8.4\times10^{28}\,\text{m}^{-3}$，
\begin{equation}
\omega_p=\sqrt{\frac{ne^2}{\varepsilon_0 m}}
=\sqrt{\frac{8.4\times10^{28}\times(1.6\times10^{-19})^2}{8.85\times10^{-12}\times9.11\times10^{-31}}}
\approx1.6\times10^{16}\,\text{rad/s}.
\end{equation}
换算为光子能量
\begin{equation}
\hbar\omega_p\approx\frac{1.055\times10^{-34}\times1.6\times10^{16}}{1.602\times10^{-19}}
\approx\boxed{10.5\,\text{eV}\approx10\,\text{eV}.}
\end{equation}

带间跃迁频率：铜呈紫红色，对绿/蓝光（$\lambda\sim500\,$nm，$\hbar\omega\sim2.5\,$eV）有较强吸收。带间跃迁（$d\to s$）阈值约为
\begin{equation}
\boxed{\hbar\omega_0\approx2\,\text{eV}.}
\end{equation}

介电函数分离实部与虚部：
\begin{align}
\varepsilon_1(\omega)&=1+\frac{\omega_p^2(\omega_0^2-\omega^2)}{(\omega_0^2-\omega^2)^2+\omega^2\tau^{-2}},\\
\varepsilon_2(\omega)&=\frac{\omega_p^2\omega\tau^{-1}}{(\omega_0^2-\omega^2)^2+\omega^2\tau^{-2}}.
\end{align}

代入 $\hbar\omega_p=10\,$eV，$\hbar\omega_0=2\,$eV，$\hbar\tau^{-1}\approx0.066\,$eV。曲线特征：
\begin{itemize}
    \item $\varepsilon_1$：在 $\omega<\omega_0$ 区域为正且较大；在 $\omega=\omega_0$ 附近出现反常色散（先升后降）；$\omega\to\infty$ 时趋于 1。
    \item $\varepsilon_2$：在 $\omega=\omega_0$ 处出现洛伦兹共振吸收峰，峰宽由 $\tau^{-1}$ 决定。
\end{itemize}

\textbf{(2) 室温下的复电导率}

由 $\varepsilon(\omega)=1+i\sigma(\omega)/(\varepsilon_0\omega)$，对比特鲁德-洛伦兹公式得
\begin{equation}
\sigma(\omega)=\frac{\varepsilon_0\omega_p^2\bigl[\omega\tau^{-1}-i(\omega_0^2-\omega^2)\bigr]}{(\omega_0^2-\omega^2)^2+\omega^2\tau^{-2}}.
\end{equation}

其实部和虚部分别为
\begin{align}
\operatorname{Re}\sigma(\omega)&=\frac{\varepsilon_0\omega_p^2\omega\tau^{-1}}{(\omega_0^2-\omega^2)^2+\omega^2\tau^{-2}},\\
\operatorname{Im}\sigma(\omega)&=-\frac{\varepsilon_0\omega_p^2(\omega_0^2-\omega^2)}{(\omega_0^2-\omega^2)^2+\omega^2\tau^{-2}}.
\end{align}

\textbf{(3) 绝对零度、纯净无缺陷时的复电导率}

$\tau\to\infty$（$\tau^{-1}\to0$），代入上式得
\begin{equation}
\boxed{\sigma(\omega)=-\frac{i\varepsilon_0\omega_p^2}{\omega_0^2-\omega^2}.}
\end{equation}
此时电导率为\textbf{纯虚数}（实部为零），意味着无耗散，响应为无损的极化振荡。
\end{derivationbox}

\begin{formulabox}[答案汇总]
\begin{center}
\begin{tabular}{ll}
\toprule
\textbf{项目} & \textbf{结果} \\
\midrule
$\omega_p$ & $\approx1.6\times10^{16}\,$rad/s，$\hbar\omega_p\approx10\,$eV \\
$\omega_0$ & $\hbar\omega_0\approx2\,$eV（铜的颜色） \\
室温 $\sigma(\omega)$ & $\displaystyle\sigma(\omega)=\frac{\varepsilon_0\omega_p^2[\omega\tau^{-1}-i(\omega_0^2-\omega^2)]}{(\omega_0^2-\omega^2)^2+\omega^2\tau^{-2}}$ \\
零温纯净 $\sigma(\omega)$ & $\displaystyle\sigma(\omega)=-\frac{i\varepsilon_0\omega_p^2}{\omega_0^2-\omega^2}$（纯虚） \\
\bottomrule
\end{tabular}
\end{center}
\end{formulabox}

\begin{insightbox}[物理图像]
\begin{itemize}
    \item 特鲁德-洛伦兹模型将自由电子的德鲁德响应（低频）与束缚电子的洛伦兹振子响应（带间跃迁）结合，成功描述了金属从红外到紫外的光学性质。
    \item 铜的红色外观源于带间跃迁对蓝绿光的吸收：当光子能量 $>2\,$eV 时，$d$ 带电子可被激发到费米面附近的 $s$ 带，导致反射率下降。
    \item 复电导率的实部对应光吸收（焦耳热），虚部对应极化响应。绝对零度下无散射，实部为零，系统表现为无损的等离子体介质。
\end{itemize}
\end{insightbox}
